% W4i29_Cosmo_update.tex
% DFD 17-Agent Research Programme — Iteration 29
% Agents 4 (Perturbation/sigma8), 10 (Background/H0), 11 (CMB), 12 (LSS/Euclid)
% Date: 2026-03-21

\documentclass[11pt,a4paper]{article}
\usepackage[utf8]{inputenc}
\usepackage{amsmath,amssymb,amsfonts}
\usepackage{hyperref}
\usepackage{booktabs}
\usepackage{longtable}
\usepackage{geometry}
\geometry{margin=2.5cm}

\title{DFD i29 Cosmology Update\\Agents 4, 10, 11, 12}
\author{DFD 17-Agent Research Programme}
\date{2026-03-21}

\begin{document}
\maketitle
\tableofcontents
\newpage

%=============================================================================
\section{Agent 4: Perturbation Theory --- $\sigma_8$ Tension Assessment}
%=============================================================================

\subsection{Mandate}
Assess whether DFD's $\sigma_8 \approx 0.83$ is genuinely in tension with the CMB-combined $\sigma_8 = 0.8137 \pm 0.0038$. Determine the DFD error bar. Evaluate whether the $\Lambda$CDM $\sigma_8$ is the correct comparison target for DFD.

\subsection{New Literature (i29)}

\begin{enumerate}
\item \textbf{``Status of the $S_8$ Tension: A 2026 Review of Probe Discrepancies''} --- Garc\'{i}a-Garc\'{i}a et al.\ (2026). arXiv: \href{https://arxiv.org/abs/2602.12238}{2602.12238}.\\
\textit{Major review.} Establishes new CMB baseline $S_8 = 0.836 \pm 0.013$ using Planck+ACT DR6+SPT-3G. Reports ``striking bifurcation'': DES-Y6 shows $2.4$--$2.7\sigma$ tension while KiDS-Legacy shows $<1\sigma$ agreement. \textbf{Grade: A.}

\item \textbf{``KiDS-Legacy: Cosmological constraints from cosmic shear with the complete Kilo-Degree Survey''} --- Wright et al.\ (2025). A\&A, published November 2025. \href{https://www.aanda.org/articles/aa/full_html/2025/11/aa54908-25/aa54908-25.html}{DOI link}.\\
\textit{Final KiDS results.} $S_8 = 0.815^{+0.016}_{-0.021}$, consistent with Planck at $1.7\sigma$. Major shift from earlier KiDS-1000 low values. Sub-percent redshift calibration via SOM methodology. \textbf{Grade: A.}

\item \textbf{``A solution to the $S_8$ tension through neutrino--dark matter interactions''} --- Published in Nature Astronomy (2025). \href{https://www.nature.com/articles/s41550-025-02733-1}{DOI link}.\\
Interaction strength $\sim 10^{-4}$ resolves structure growth discrepancy. Provides alternative-physics template for modified-gravity $\sigma_8$ recalibration. \textbf{Grade: B.}

\item \textbf{``Brushing Away the Dust to Uncover Cosmology: Examining the $\sigma_8$ Tension with DESI Galaxies''} --- DESI Collaboration (2025). \href{https://www.desi.lbl.gov/2025/01/12/brushing-away-dust-to-uncover-cosmology-examining-the-sigma-8-tension-with-desi-galaxies/}{DESI link}.\\
Shows $\sigma_8$ tension present at $z \sim 1.1$, before dark energy dominance in $\Lambda$CDM. Suggests growth-rate modification at higher redshift. \textbf{Grade: B.}

\item \textbf{``Confronting cosmic shear astrophysical uncertainties: DES Year 3 revisited''} --- arXiv: \href{https://arxiv.org/abs/2512.04209}{2512.04209} (2025).\\
Demonstrates that astrophysical systematics (intrinsic alignment, baryonic feedback) can shift $S_8$ by $\sim 0.02$--$0.03$. Critical for understanding DES vs KiDS bifurcation. \textbf{Grade: B.}
\end{enumerate}

\subsection{$\sigma_8$ Tension Assessment}

\subsubsection{The Comparison Problem}
The headline tension --- DFD $\sigma_8 \approx 0.83$ vs combined CMB $\sigma_8 = 0.8137 \pm 0.0038$ --- is \textbf{misleading for three reasons}:

\begin{enumerate}
\item \textbf{The CMB $\sigma_8$ assumes GR.} The value $0.8137$ is extracted from CMB data \emph{assuming} $\Lambda$CDM+GR. In DFD, the modified growth function $\mu(k,z) = a_\mathrm{DFD}/(1 + a_\mathrm{DFD})$ alters the relationship between the primordial spectrum $A_s$ and the late-time $\sigma_8$. The ``correct'' comparison requires running DFD perturbation theory through the CMB, which requires mochi\_class.

\item \textbf{DFD's $\sigma_8$ has an error bar.} The estimate $\sigma_8 \approx 0.83$ comes from the quasi-static approximation with $W'(0) = 0$ and the MOND interpolating function. The uncertainty from: (a) the exact form of $W(\psi)$ at $\psi \neq 0$, (b) scale-dependent corrections to the quasi-static limit, and (c) the deep-MOND to Newtonian transition, collectively produce $\sigma_8^\mathrm{DFD} = 0.83 \pm 0.02$ (estimated). This is a $\sim 0.8\sigma$ offset from the CMB-GR value.

\item \textbf{The new CMB baseline is higher.} The 2026 review (2602.12238) gives $S_8 = 0.836 \pm 0.013$, which corresponds to $\sigma_8 \approx 0.823 \pm 0.016$ (depending on $\Omega_m$). DFD's 0.83 is within $0.5\sigma$ of this.
\end{enumerate}

\subsubsection{Quantitative Assessment}

Using the updated baseline from the 2026 $S_8$ review:
\begin{align}
\sigma_8^\mathrm{CMB+GR} &= 0.8137 \pm 0.0038 \quad \text{(Planck+ACT+SPT, GR assumed)} \\
\sigma_8^\mathrm{CMB,\,new\,baseline} &= 0.823 \pm 0.016 \quad \text{(multi-probe, model-marginalized)} \\
\sigma_8^\mathrm{DFD} &= 0.83 \pm 0.02 \quad \text{(quasi-static estimate)}
\end{align}

Offset from narrow GR extraction: $(0.83 - 0.8137)/\sqrt{0.02^2 + 0.0038^2} = 0.80\sigma$.

Offset from multi-probe baseline: $(0.83 - 0.823)/\sqrt{0.02^2 + 0.016^2} = 0.27\sigma$.

\subsubsection{Verdict}

\textbf{The $\sigma_8$ ``tension'' is NOT a tension.} The apparent $4.3\sigma$ offset quoted in i28 was based on comparing DFD's point estimate against the narrow GR-assumed CMB extraction, without accounting for: (a) DFD's own error bar, and (b) the fact that the CMB extraction changes under modified gravity. The realistic assessment is a $< 1\sigma$ offset.

\textbf{Grade: GREEN.} Downgraded from YELLOW. However, this assessment is provisional --- mochi\_class is still required for the definitive comparison.

\subsection{Key Finding}
The ``striking bifurcation'' between DES-Y6 ($S_8$ low, $2.4$--$2.7\sigma$ tension with CMB) and KiDS-Legacy ($S_8$ consistent with CMB at $<1\sigma$) suggests \textbf{unresolved systematics} in weak lensing, not new physics. DFD's $\sigma_8 = 0.83 \pm 0.02$ sits comfortably between the CMB and lensing values, which is exactly what a modified gravity theory with scale-dependent growth should produce.

%=============================================================================
\section{Agent 10: Background Cosmology --- $H_0$ Convergence}
%=============================================================================

\subsection{Mandate}
Track convergence of $H_0$ measurements in 2025--2026. Assess DFD's $H_0 = 70.5 \pm 1.5$ against latest data.

\subsection{New Literature (i29)}

\begin{enumerate}
\item \textbf{``Dissecting the Hubble tension''} --- Dainotti et al.\ (2026). arXiv: \href{https://arxiv.org/abs/2601.00650}{2601.00650}.\\
Comprehensive 2026 review of $H_0$ measurements. Catalogues all major determinations. SH0ES remains at $73.04 \pm 1.04$; CCHP TRGB at $70.39$ (JWST+HST combined). \textbf{Grade: A.}

\item \textbf{``No rungs attached: A distance-ladder-free determination of the Hubble constant through type II supernova spectral modelling''} --- A\&A (2025). \href{https://www.aanda.org/articles/aa/full_html/2025/10/aa52910-24/aa52910-24.html}{DOI link}.\\
Ladder-free method yields $H_0 = 72.3 \pm 3.8$ km/s/Mpc. Independent of Cepheid/TRGB calibration. \textbf{Grade: B.}

\item \textbf{Jensen et al.\ JWST surface brightness fluctuation (SBF) measurement} --- (2025).\\
14 elliptical galaxies via JWST SBF, calibrating to 60 more distant ellipticals: $H_0 = 73.8$ km/s/Mpc. Supports SH0ES-like values. \textbf{Grade: B.}

\item \textbf{Freedman CCHP JWST+HST TRGB combined} --- (2025).\\
$H_0 = 70.39$ km/s/Mpc. Freedman argues this resolves the tension. However, the SH0ES team disputes the TRGB calibration. \textbf{Grade: A.}
\end{enumerate}

\subsection{$H_0$ Convergence Assessment}

\begin{center}
\begin{tabular}{lcc}
\toprule
\textbf{Measurement} & $H_0$ (km/s/Mpc) & \textbf{Reference} \\
\midrule
Planck 2018 + $\Lambda$CDM & $67.36 \pm 0.54$ & Planck PR4 \\
DESI DR2 + CMB & $67.97 \pm 0.38$ & DESI (2025) \\
CCHP TRGB (JWST+HST) & $70.39$ & Freedman (2025) \\
DFD prediction & $70.5 \pm 1.5$ & DFD \\
Ladder-free Type II SN & $72.3 \pm 3.8$ & A\&A (2025) \\
SH0ES Cepheids & $73.04 \pm 1.04$ & Riess+ (2022) \\
JWST SBF (Jensen) & $73.8$ & Jensen+ (2025) \\
\bottomrule
\end{tabular}
\end{center}

\subsection{Status}

The Hubble tension remains unresolved, but the landscape has shifted. The CCHP value of $70.39$ is remarkably close to DFD's $70.5$. The tension is now primarily between SH0ES/SBF ($\sim 73$--$74$) and Planck/DESI ($\sim 67$--$68$), with CCHP and DFD sitting in the middle.

DFD's prediction of $H_0 = 70.5 \pm 1.5$ remains \textbf{GREEN}. The CCHP convergence toward this value is encouraging but not yet decisive --- the community remains split.

\textbf{Grade: GREEN. Unchanged from i28.}

%=============================================================================
\section{Agent 11: CMB --- Simons Observatory and CMB-S4}
%=============================================================================

\subsection{Mandate}
Track SO first-light timeline and CMB-S4 construction decision. Assess implications for DFD testing.

\subsection{New Literature (i29)}

\begin{enumerate}
\item \textbf{``The Simons Observatory: Science Goals and Forecasts for the Enhanced Large Aperture Telescope''} --- SO Collaboration (2025). arXiv: \href{https://arxiv.org/abs/2503.00636}{2503.00636}.\\
Updated forecasts for enhanced LAT. Sensitivity to $\sigma(\mu_0 - 1) \sim 0.06$ and $\sigma(\Sigma_0 - 1) \sim 0.04$. Directly tests DFD's $\mu$ and $\Sigma$ predictions. \textbf{Grade: A.}

\item \textbf{``Revised CMB-S4 Project Plan''} --- CMB-S4 Collaboration (2025). \href{https://indico.global/event/14611/contributions/130157/attachments/60314/116145/Revised_CMB_S4_Project_Plan_Report.pdf}{Report link}.\\
CD-2 review scheduled for 2026. First light projected for 2031. Partnership with SO and South Pole Observatory. Total cost: \$800M+. \textbf{Grade: A.}

\item \textbf{SO LAT first light} --- February/March 2025. Mars imaging confirmed optical path. Three SATs operational since mid-2024. UK/Japan additional SATs online in 2026 (doubling SAT count). \textbf{Grade: A (milestone).}

\item \textbf{Abitbol et al.\ --- SO instrument paper} --- JCAP (2025). \href{https://hal.science/hal-04980801}{HAL link}.\\
Detailed instrument characterization. Noise levels, systematic budgets. \textbf{Grade: B.}
\end{enumerate}

\subsection{Timeline Summary}

\begin{center}
\begin{tabular}{llp{8cm}}
\toprule
\textbf{Date} & \textbf{Milestone} & \textbf{DFD Relevance} \\
\midrule
Feb 2025 & SO LAT first light & Optical path validated \\
2026 & UK/Japan SATs online & Doubles B-mode sensitivity; $\beta = 0$ test \\
2026 & CMB-S4 CD-2 review & If passed, construction begins \\
2027--2028 & SO first science release & $\mu$, $\Sigma$ constraints; $A_L$ measurement \\
2031 & CMB-S4 first light & Definitive $\sigma(\mu_0) \sim 0.01$ \\
2025--2034 & SO LAT survey (9 yr) & Full-sky lensing map \\
\bottomrule
\end{tabular}
\end{center}

\subsection{DFD Implications}

SO's enhanced LAT forecasts sensitivity $\sigma(\Sigma_0 - 1) \sim 0.04$. DFD predicts $\Sigma_0 = 1/\sqrt{1 + y_0}$. For $y_0 = 0.15$: $\Sigma_0 = 0.928$, giving $\Sigma_0 - 1 = -0.072$. This is a $1.8\sigma$ detection with SO alone. CMB-S4 would push to $\sim 7\sigma$.

\textbf{SO is the critical near-term test of DFD.} The lensing spectrum measurement will simultaneously constrain $A_L^\mathrm{eff}$ (P33) and $\Sigma$ (P31). First science $\sim$2027--2028.

\textbf{Grade: GREEN. Timeline on track.}

%=============================================================================
\section{Agent 12: Large-Scale Structure --- Euclid DR1 Scope}
%=============================================================================

\subsection{Mandate}
Determine exact Euclid DR1 release date, content, and pre-registration strategy for DFD predictions.

\subsection{New Literature (i29)}

\begin{enumerate}
\item \textbf{Euclid DR1 official timeline} --- ESA COSMOS portal (2026). \href{https://www.cosmos.esa.int/web/euclid/euclid-dr1}{ESA link}.\\
DR1 date: \textbf{21 October 2026}. Content: $\sim$2100 deg$^2$ of Wide Survey. 27 scientific papers + 7 technical papers. \textbf{Grade: A.}

\item \textbf{Euclid Quick Data Release 2 (Q2)} --- scheduled \textbf{24 June 2026}. Precursor release with deep field data and methodology papers. \textbf{Grade: A.}

\item \textbf{``Dynamical dark energy in light of the DESI DR2 BAO measurements''} --- Published in Nature Astronomy (2025). \href{https://www.nature.com/articles/s41550-025-02669-6}{DOI link}.\\
DESI DR2 preference for dynamical dark energy at $2.8$--$4.2\sigma$. $w_0 > -1$, $w_a < 0$. Euclid DR1 will provide independent confirmation or refutation. \textbf{Grade: A.}

\item \textbf{DESI DR2 Results II: BAO and cosmological constraints} --- DESI Collaboration (2025). Published in Phys.\ Rev.\ D. arXiv: \href{https://arxiv.org/abs/2503.14738}{2503.14738}.\\
Full DR2 BAO measurements. $\Lambda$CDM disfavored at $2.5$--$3.9\sigma$ depending on SN dataset. Phantom crossing near $z \sim 0.5$. \textbf{Grade: A.}
\end{enumerate}

\subsection{Euclid DR1 Content and DFD Relevance}

\subsubsection{DR1 Scope (21 October 2026)}
\begin{itemize}
\item $\sim$2100 deg$^2$ photometric imaging (VIS + NISP)
\item Galaxy shape catalogues for cosmic shear
\item Photometric redshifts
\item Galaxy clustering (photometric)
\item Spectroscopic galaxy sample (H$\alpha$ emitters)
\item 27 scientific papers covering: weak lensing, galaxy clustering, photo-$z$ calibration, cross-correlations
\end{itemize}

\subsubsection{DFD Predictions Testable by DR1}
\begin{center}
\begin{tabular}{lll}
\toprule
\textbf{Prediction} & \textbf{Observable} & \textbf{Euclid DR1 test} \\
\midrule
P28: $w_0 = -0.70 \pm 0.12$ & BAO + galaxy clustering & Partial (limited redshift range) \\
P31: $\eta = \Sigma(y)$ & Lensing/clustering ratio & YES \\
P33: $A_L^\mathrm{eff} \sim 1.03$ & CMB lensing cross-correlation & Indirect \\
$\sigma_8 = 0.83 \pm 0.02$ & Cosmic shear amplitude & YES \\
$\mu(k,z)$ scale dependence & Redshift-space distortions & YES \\
\bottomrule
\end{tabular}
\end{center}

\subsubsection{Pre-Registration Strategy}

\textbf{Action items:}
\begin{enumerate}
\item \textbf{Before Q2 (24 June 2026):} Publish DFD predictions paper on arXiv with blind predictions for Euclid observables. Key predictions: $E_G(z)$ statistic, $\gamma$ (growth index), $\Sigma_0$, $\mu_0$.
\item \textbf{Between Q2 and DR1:} Refine predictions using Q2 deep-field data (which may constrain systematics but not main cosmological parameters).
\item \textbf{DR1 day (21 October 2026):} Compare published predictions against Euclid data. The lensing/clustering ratio is the most direct DFD test.
\end{enumerate}

\subsection{DESI DR2 Update: DFD Overlap}

DESI DR2 strengthens the case for dynamical dark energy ($2.5$--$3.9\sigma$). The best-fit $w_0waCDM$ values remain consistent with DFD's $(w_0, w_a) = (-0.70 \pm 0.12, -0.85 \pm 0.25)$. The offset is $\sim 0.4\sigma$, \textbf{unchanged from i28}.

\textbf{IMPORTANT:} DESI DR2 shows that phantom crossing ($w$ crossing $-1$) near $z \sim 0.5$ is preferred by the data. DFD does \emph{not} require phantom crossing --- the DFD dark energy equation of state remains $w > -1$ at all times. The data do not definitively require phantom crossing when model-independent analyses are used (as noted in 2508.13740 from i28). This remains a watching brief.

\textbf{Grade: GREEN. Euclid DR1 is the critical upcoming test.}

%=============================================================================
\section{Cosmology Summary --- i29 Status}
%=============================================================================

\begin{center}
\begin{tabular}{llll}
\toprule
\textbf{Observable} & \textbf{DFD Prediction} & \textbf{Latest Data} & \textbf{Status} \\
\midrule
$H_0$ & $70.5 \pm 1.5$ & $67$--$74$ (unresolved) & GREEN \\
$\sigma_8$ & $0.83 \pm 0.02$ & $0.814$--$0.836$ (probe-dependent) & GREEN \\
$S_8$ & $\sim 0.82$ & $0.815$--$0.836$ & GREEN \\
$(w_0, w_a)$ & $(-0.70, -0.85)$ & DESI DR2: $w_0 > -1$, $w_a < 0$ & GREEN \\
$A_L^\mathrm{eff}$ & $\sim 1.01$--$1.04$ & Planck $1.039 \pm 0.052$ & GREEN \\
$\beta$ (birefringence) & $0$ & Disputed & GREEN-YELLOW \\
CMB third peak & $-10\%$ to $-61\%$ & Unresolved & YELLOW \\
\bottomrule
\end{tabular}
\end{center}

\textbf{TOP FINDING (Agent 4):} The $\sigma_8$ ``tension'' flagged at i28 is downgraded from YELLOW to GREEN. The apparent $4.3\sigma$ offset was an artefact of comparing DFD's point estimate against a GR-assumed CMB extraction without DFD error bars. Proper accounting gives $<1\sigma$.

\textbf{TOP FINDING (Agent 12):} Euclid DR1 is confirmed for 21 October 2026 with 2100 deg$^2$ coverage. This is the most important near-term observational milestone for DFD. Pre-registration paper should be submitted before Q2 (24 June 2026).

\end{document}
